\documentclass[12pt, a4paper]{article}

% Packages
\usepackage[utf8]{inputenc}
\usepackage{amsmath, amssymb, amsthm}
\usepackage{geometry}
\usepackage{hyperref}
\usepackage{titlesec}
\usepackage{mathrsfs}
\usepackage{booktabs}
\usepackage{array}

% Theorem environments
\newtheoremstyle{note}{3pt}{3pt}{}{}{\bfseries}{.}{ }{}
\theoremstyle{note}
\newtheorem*{note}{Note}

\newtheorem{theorem}{Theorem}[section]
\newtheorem{lemma}[theorem]{Lemma}
\newtheorem{corollary}[theorem]{Corollary}
\newtheorem{proposition}[theorem]{Proposition}
\newtheorem{definition}[theorem]{Definition}
\newtheorem{remark}{Remark}[section]
\newtheorem{example}[theorem]{Example}

% Page geometry
\geometry{margin=1in}

% Custom commands
\newcommand{\floor}[1]{\left\lfloor #1 \right\rfloor}
\newcommand{\Real}{\mathbb{R}}
\newcommand{\Complex}{\mathbb{C}}
\newcommand{\Int}{\mathbb{Z}}

% Hyperlinks
\hypersetup{
    colorlinks=true,
    linkcolor=blue,
    filecolor=magenta,
    urlcolor=cyan,
}

% Title information
\title{\textbf{An Alternative Analytic Representation of the Floor Function and Related Discrete Operations}}
\author{Mohammad Rakibul Islam \\ \textit{High School Student} \\ \texttt{rakibulofficial1017@gmail.com}}
\date{\today}

\begin{document}

\maketitle

\begin{abstract}
The floor function $\floor{x}$ is traditionally defined through piecewise cases, infinite series, or auxiliary discrete operators. Existing analytic representations like $x - \frac{1}{2} + \pi^{-1}\arctan(\cot(\pi x))$ exhibit vertical asymptotes at integers, complicating symbolic differentiation and complex-domain analysis.

This paper presents an alternative closed-form expression for $\floor{x}$ constructed from elementary trigonometric functions, derivatives, and arithmetic operations. Using $g(x) = \pi^{-1}\arccos(\cos(\pi x))$ with one-sided limits $\lim_{n \to x^+} g'(n)$ to handle integer boundaries, this construction avoids piecewise definitions and auxiliary functions such as $\operatorname{sgn}$. Unlike prior formulations using cotangent, this approach mitigates unbounded singularities at the cost of requiring numerical evaluation of derivative signs.

For complex inputs $z = a + bi$, empirical evaluation reveals system-dependent behavior between computational platforms. From this foundation, explicit formulas for ceiling, rounding, modulo, and precision rounding are derived. A generalized block-extraction operator $E_{\text{right}}(x, y, b)$ anchored at the decimal point enables direct element access using base-10 encoding. A two-function base conversion framework handles arbitrary radix transformations for bases $2 \leq B \leq 10$. Analytic data structures demonstrate how numeric arrays can be encoded as single integers with recovery—technique related to Gödel numbering. Logical predicates (equality, order comparisons) and number-theoretic functions (divisor counting $\nu(x)$, primality testing) are constructed analytically. Finally, a Pythagorean triplet parameterization valid for complex arguments is derived. Numerical verification achieves high success rates across standard test cases; IEEE 754 floating-point implementations document precision limitations near boundaries and for large numbers as expected. These results illustrate connections between discrete computational logic and continuous calculus, providing alternative tools for symbolic computation.

\noindent\textbf{Keywords:} floor function, analytic representation, discrete mathematics, numerical methods, symbolic computation, complex analysis, base conversion, data encoding

\end{abstract}

\tableofcontents
\newpage

% ============================================================
\section{Introduction}
% ============================================================

\subsection{The Discrete-Continuous Divide}

The floor function $\floor{x}$, returning the greatest integer less than or equal to $x$, is fundamental in number theory, combinatorics, and computer science. Despite its ubiquity, representing it as a single analytic expression poses challenges. Standard treatments define it via piecewise cases:
\[
\floor{x} = n \quad \text{if } n \leq x < n+1, \quad n \in \Int,
\]
or through infinite Fourier series.

\subsection{Motivation}

Existing approaches face specific obstacles:
\begin{enumerate}
    \item \textbf{Piecewise definitions}: Incompatible with symbolic differentiation and integration in many computer algebra systems.
    \item \textbf{Fourier series}: Require infinite summation and converge slowly near integers.
    \item \textbf{Known analytic formulas}: Representations like $x - \frac{1}{2} + \pi^{-1}\arctan(\cot(\pi x))$ possess vertical asymptotes at every integer $x \in \Int$, preventing stable symbolic manipulation.
    \item \textbf{Calculator incompatibility}: Standard calculators lack built-in floor functions, motivating the search for pure arithmetic-trigonometric formulations.
\end{enumerate}

This work explores an alternative approach using a derivative-based switching mechanism that trades singularity problems for numerical sampling requirements.

\subsection{Prior Work}

A well-known analytic representation exists (Mathematics Stack Exchange, various references):
\[
\lfloor x \rfloor = x - \frac{1}{2} + \frac{1}{\pi}\arctan(\cot(\pi x)).
\]
However, this formula exhibits vertical asymptotes at every integer due to the $\cot(\pi x)$ term. Its complex extension grows exponentially along the imaginary axis due to the hyperbolic nature of $\cot(\pi z)$.

Additionally, techniques for encoding sequences as single integers trace back to Gödel's incompleteness theorem work (1931), and indicator functions using floor operations are well-established in mathematical logic.

\subsection{Contributions}

This paper presents:
\begin{itemize}
    \item An alternative closed-form expression for $\floor{x}$ using arccos/cos, derivatives, and arithmetic, with discussion of singularity versus numerical evaluation trade-offs.
    \item Documentation of system-dependent behavior for complex inputs across computational platforms.
    \item Algebraic derivations of ceiling, rounding, modulo, and precision rounding.
    \item A generalized block-extraction operator for variable-width element access, with both left-to-right and right-to-left variants documented.
    \item A base conversion framework for bases $2 \leq B \leq 10$.
    \item Data encoding schemes related to Gödel numbering.
    \item Logical predicates and number-theoretic functions constructed analytically.
    \item A Pythagorean triplet parameterization extending to complex arguments.
    \item Comprehensive computational verification (250/251 test cases passed).
\end{itemize}

\textbf{Disclaimer}: This paper aims to be intellectually honest about the relationship between this representation and existing approaches. The goal is not to claim priority but to explore an alternative formulation.

% ============================================================
\section{The Alternative Floor Function Formulation}
\label{sec:floor}
% ============================================================

\subsection{Mathematical Formulation}

\begin{definition}\label{def:g}
Define the core oscillatory function:
\[
g(x) = \frac{1}{\pi}\arccos\!\left(\cos\!\left(\pi x\right)\right).
\]
This produces a triangular waveform with period 2, ranging continuously from $[0,1]$ across each unit interval.
\end{definition}

\begin{theorem}[Alternative Floor Representation]\label{thm:floor}
For any real number $x$, the floor function may be represented as:
\[
F(x) = x - g\!\left(x + \frac{g'^{+}(x) - 1}{2}\right),
\]
where $g'^{+}(x) = \lim_{n \to x^+} g'(n)$ denotes the one-sided derivative limit approaching $x$ from the positive direction. Then $F(x) = \floor{x}$ for all $x \in \Real$ where the derivative exists or the limit is properly defined.
\end{theorem}

\begin{proof}
The key insight is that $g'(x)$ alternates between $+1$ and $-1$ depending on whether $x$ lies in even-numbered intervals $[2k, 2k+1)$ or odd-numbered intervals $[2k+1, 2k+2)$. The correction term $\frac{g'^{+}(x) - 1}{2}$ evaluates to $0$ when $g'^{+}(x) = +1$ (even intervals) and $-1$ when $g'^{+}(x) = -1$ (odd intervals), shifting inputs appropriately so that $g(\cdot)$ returns the correct fractional offset.

Specifically:
\begin{itemize}
    \item For $x \in [2k, 2k+1)$: $g'^{+}(x) = +1$, so correction $= 0$, and $F(x) = x - g(x)$. Since $g(x) \in [0,1)$ on this interval with $g(x) = x - 2k$, we get $F(x) = x - (x - 2k) = 2k = \floor{x}$.
    \item For $x \in [2k+1, 2k+2)$: $g'^{+}(x) = -1$, so correction $= -1$, and $F(x) = x - g(x - 1)$. Since $g(x-1) = (x - 1) - 2k = x - 2k - 1$, we get $F(x) = x - (x - 2k - 1) = 2k + 1 = \floor{x}$.
\end{itemize}

At exact integers $x = n$, the one-sided limit $g'^{+}(n)$ ensures evaluation from the right-hand side.
\end{proof}

\subsection{The Switching Mechanism}

\begin{lemma}[Derivative Switch]\label{lem:switch}
For any non-integer $x \in \Real \setminus \Int$, the derivative $g'(x)$ satisfies:
\[
g'(x) = 
\begin{cases} 
+1 & \text{if } x \in [2k, 2k+1) \\
-1 & \text{if } x \in [2k+1, 2k+2)
\end{cases} \text{ for } k \in \Int.
\]
\end{lemma}

\begin{proof}
By definition $g(x) = \pi^{-1}\arccos(\cos(\pi x))$. On intervals where $\cos(\pi x)$ is monotonically decreasing, the chain rule yields $g'(x) = +1$; where it is increasing, $g'(x) = -1$. The sign flips at each integer where $\cos(\pi x)$ reaches its extrema.
\end{proof}

\subsection{Integer Boundary Resolution}

\begin{remark}[Boundary Handling via One-Sided Limits]\label{rem:boundary}
At exact integers $x = n \in \Int$, $g'(n)$ is formally undefined due to corner points in the triangular waveform. We resolve this by evaluating $\lim_{n \to x^+} g'(n)$, ensuring consistent evaluation from the right-hand side without explicit piecewise definitions.

In numerical implementation, this limit is approximated by sampling $g$ at $x + \epsilon$ and $x + 2\epsilon$ for small $\epsilon > 0$ (empirically $\epsilon = 10^{-7}$) and determining whether $g$ is increasing or decreasing.

\textbf{Important Caveat}: This numerical sampling step is itself a discrete operation, albeit implicit. A reviewer might argue this trades explicit piecewise definition for implicit numerical discretization. Both formulations solve the same underlying problem: determining which interval contains the input.
\end{remark}

\begin{remark}[Adaptive Perturbation Scaling]\label{rem:adaptive_eps}
For large numbers exceeding $\sim 10^{11}$, fixed perturbation $\epsilon = 10^{-7}$ falls below IEEE 754 double-precision machine epsilon at that scale. Resolution requires adaptive scaling: $\epsilon = \max(10^{-7}, |x| \cdot 10^{-10})$, clamped to prevent excessive perturbation. All subsequent tests employ adaptive perturbation.
\end{remark}

\subsection{Extension to the Complex Plane}

The function $F(x)$ can be evaluated for complex arguments via analytic continuation of the constituent functions. Unlike the $\cot$-based formula, the arccos/cos formulation remains bounded for finite inputs in certain regimes.

\begin{table}[h!]
\centering
\caption{Complex Domain Behavior by Computational System}
\begin{tabular}{l|c|c}
Input $z$ & Python (cmath) & Wolfram Alpha \\ \hline
$1 + i$ & $1$ & $1$ \\
$-1 + i$ & $-1$ & $-1$ \\
$1 - i$ & $1 - 2i$ & $1$ \\
$-1 - i$ & $-1 + 2i$ & $-1$ \\
$2 - 3i$ & $2 + 6i$ & $2$ \\
\end{tabular}
\end{table}

\begin{remark}[Observed System Divergence]\label{rem:system_divergence}
Two distinct behaviors are observed depending on the computational platform:

\textbf{Python (IEEE 754, cmath module):}
\[
F_{\text{Python}}(a + bi) = 
\begin{cases} 
\lfloor a \rfloor & \text{if } b \geq 0 \\
\lfloor a \rfloor + 2bi & \text{if } b < 0
\end{cases}
\]

\textbf{Wolfram Alpha (symbolic computation):}
\[
F_{\text{WA}}(a + bi) = \lfloor a \rfloor \quad \text{for all } b \in \Real
\]

This discrepancy suggests that the residue behavior is \textbf{implementation-dependent}, arising from how each system evaluates $\arccos(\cos(\pi z))$ across complex branch cuts. The Python implementation follows direct IEEE 754 floating-point evaluation, while Wolfram Alpha may apply symbolic simplification rules that suppress the imaginary component.
\end{remark}

\begin{remark}[Scope Limitation]\label{rem:complex_limitation}
Given this unresolved discrepancy, the primary contributions of this paper concern the \textbf{real domain} $x \in \Real$, where the formula achieves high accuracy across all test cases. Complex-domain extensions are documented here as \textbf{preliminary observations} pending rigorous analysis. Users should verify complex behavior with their specific computational environment.
\end{remark}

% ============================================================
\section{Derived Discrete Operations}
\label{sec:derived}
% ============================================================

Since $F(x)$ represents the floor function for real inputs, standard discrete operations can be derived algebraically.

\begin{corollary}[Derived Discrete Operators]\label{cor:derived}
The following operations are expressible in terms of $F(x)$:
\begin{align}
\lceil x \rceil &= -F(-x), \\
\operatorname{round}(x) &= F(x + 0.5), \\
R(x, P) &= \frac{F(x \cdot 10^{P}+0.5)}{10^{P}} \quad \text{(Precision Rounding)}, \\
x \bmod y &= x - y \cdot F\!\left(\frac{x}{y}\right) \quad \text{(Remainder)}.
\end{align}
Indicator functions for divisibility testing (returns 1 for true, 0 for false):
\begin{align}
\delta_{\operatorname{div}}(x, y) &= x - y \cdot F\!\left(-\frac{\left(x-y\cdot F\!\left(\frac{x}{y}\right)\right)^{2}}{\left(x-y\cdot F\!\left(\frac{x}{y}\right)\right)^{2}+1}\right)+1 \quad \text{(Divisibility testing)}, \\
\delta_{\operatorname{whole}}(x) &= F\!\left(F(x)-x\right)+1 \quad \text{(Whole number check)}.
\end{align}
\end{corollary}

\begin{proof}
All results follow directly from properties of the floor function: symmetry gives $\lceil x \rceil = -\lfloor -x \rfloor$, nearest-integer rounding adds $0.5$ before flooring, precision scaling adjusts decimal positions, and Euclidean division defines remainder as $x - y\lfloor x/y \rfloor$. For the indicator functions: when $y \mid x$, the remainder $(x - y\cdot F(x/y))$ vanishes, making the squared fraction $0$, so $F(0) = 0$ and the expression returns $1$. When $y \nmid x$, the remainder is nonzero, the fraction lies in $(-1, 0)$, and $F$ of this quantity is $-1$, yielding $0$. Similarly, $\delta_{\operatorname{whole}}$ checks whether $F(x) - x = 0$, which occurs precisely when $x$ is an integer.
\end{proof}

% ============================================================
\section{Analytic Base Conversion and Digit Extraction}
% ============================================================

\subsection{Left-to-Right Block Extraction (Original Formulation)}

\begin{theorem}[Left-to-Right Block Extraction]\label{thm:left_block}
For a positive number $x$, extraction index $y \geq 1$ (counting from left), and block size $b \geq 1$:
\[
E(x, y, b) = F\!\left(x \cdot 10^{b-1-F\left(\log_{10}(x)-by+b\right)}\right) - F\!\left(x \cdot 10^{by-b-F\left(\log_{10}(x)+1\right)}\right) \cdot 10^{b}.
\]
Then $E(x, y, b)$ extracts the $y$-th block of size $b$ from the left side of $x$ (most significant digits first).
\end{theorem}

\begin{remark}[Motivation for Right-to-Left Revision]
\label{rem:left_motivation}
While $E(x, y, b)$ correctly extracts blocks from left to right, it creates an indexing mismatch when applied to array encoding. Since base-10 positional notation stores element $a_0$ at the least significant (rightmost) digits, the left-to-right operator reverses array indexing: retrieving $A[n]$ requires computing $E(N,\; L-n,\; b)$, where $L$ is the array length. This creates two problems:
\begin{enumerate}
    \item \textbf{Index dependency on length}: Every retrieval depends on knowing $L$ in advance.
    \item \textbf{Index instability}: Appending an element changes $L$, shifting all indices.
\end{enumerate}
These issues motivated the development of the right-to-left extraction operator $E_{\text{right}}$, which aligns naturally with positional notation and enables direct indexing $A[n] = E_{\text{right}}(N, n+1, d)$.
\end{remark}

\subsection{Right-to-Left Block Extraction (Revised Formulation)}
\label{sec:block}

\begin{theorem}[Right-to-Left Block Extraction]\label{thm:block}
For a positive number $x$, extraction index $y \geq 1$ (counting from decimal point), and block size $b \geq 1$:
\[
E_{\text{right}}(x, y, b) = F\!\left(\frac{x - F\!\left(\frac{x}{10^{by}}\right) \cdot 10^{by}}{10^{b(y-1)}}\right).
\]
Then $E_{\text{right}}(x, y, b)$ extracts the $y$-th block of size $b$ from the right side of $x$ (least significant digits first), anchored at the decimal point.
\end{theorem}

\begin{remark}[Decimal Point Anchoring]\label{rem:decimal_anchor}
\textbf{Critical Distinction}: Unlike naive digit extraction which counts from the absolute rightmost digit, $E_{\text{right}}$ anchors at the decimal point. This means:
\begin{itemize}
    \item $E_{\text{right}}(12345.678, 1, 1) = 5$ (not 8)
    \item $E_{\text{right}}(12345.678, 2, 1) = 4$
    \item $E_{\text{right}}(12345.678, 3, 1) = 3$
\end{itemize}
Fractional digits (right of the decimal) require separate handling via scaling.
\end{remark}

\begin{remark}[Comparison of Approaches]\label{rem:comparison}
Both operators achieve the same goal but differ in practical characteristics:

\begin{table}[ht!]
\centering
\caption{Left vs.\ Right Block Extraction Comparison}
\begin{tabular}{l|c|c}
Feature & $E(x,y,b)$ (Left-to-Right) & $E_{\text{right}}(x,y,b)$ (Right-to-Left) \\ \hline
Formula complexity & Higher (nested exponents) & Lower (simple powers) \\
Dependency on $\log(x)$ & Required & Not needed \\
Computational cost & Higher & Lower \\
Natural ordering & Matches reading order & Matches positional notation \\
Array indexing & Reversed: $A[n] = E(N, L{-}n, d)$ & Direct: $A[n] = E_{\text{right}}(N, n{+}1, d)$ \\
Length dependency & Requires knowing $L$ a priori & No length dependency \\
Stability on append & Indices shift & Indices unchanged \\
\end{tabular}
\end{table}

The left-to-right formula is retained for completeness and may be preferable in applications where natural reading order matters. However, all subsequent data structure constructions in this paper use $E_{\text{right}}$.
\end{remark}

\subsection{Base Conversion Framework}

\begin{theorem}[Decimal-to-Target Base Conversion]\label{thm:encode}
For a positive integer $N$ in decimal representation and target base $B_{\text{to}} \geq 2$:
\[
O(N, B_{\text{to}}) = \sum_{n=0}^{F(\log_{B_{\text{to}}}(N))}\left(B_{\text{to}} \cdot \left(\frac{F\!\left(\frac{N}{B_{\text{to}}^{n}}\right)}{B_{\text{to}}} - F\!\left(\frac{F\!\left(\frac{N}{B_{\text{to}}^{n}}\right)}{B_{\text{to}}}\right)\right) \cdot 10^{n}\right).
\]
Then $O(N, B_{\text{to}})$ returns a number whose decimal digits are the base-$B_{\text{to}}$ digits of $N$.
\end{theorem}

\begin{example}
$O(13, 2)$: Each summand extracts one binary digit and places it at the corresponding decimal position.
\begin{align*}
n=0: &\quad F(13/1) = 13, \quad 13/2 = 6.5, \quad F(6.5) = 6, \quad 2(6.5 - 6) = 1, \quad \cdot 10^0 = 1 \\
n=1: &\quad F(13/2) = 6, \quad 6/2 = 3, \quad F(3) = 3, \quad 2(3 - 3) = 0, \quad \cdot 10^1 = 0 \\
n=2: &\quad F(13/4) = 3, \quad 3/2 = 1.5, \quad F(1.5) = 1, \quad 2(1.5 - 1) = 1, \quad \cdot 10^2 = 100 \\
n=3: &\quad F(13/8) = 1, \quad 1/2 = 0.5, \quad F(0.5) = 0, \quad 2(0.5 - 0) = 1, \quad \cdot 10^3 = 1000
\end{align*}
Sum: $1 + 0 + 100 + 1000 = 1101$. Thus $O(13, 2) = 1101$.
\end{example}

\subsubsection{Source-to-Decimal Conversion: $I(N, B_{\text{from}})$}

\begin{definition}[Base Decoding]\label{def:decode}
Given a number $N$ whose decimal digits represent a value in base $B_{\text{from}}$, decode it as:
\[
I(N, B_{\text{from}}) = \sum_{n=1}^{F(\log_{10}(N) + 1)} \left(E_{\text{right}}(N, n, 1) \cdot B_{\text{from}}^{(n-1)}\right).
\]
Each decimal digit is treated as a coefficient of $B_{\text{from}}$ raised to its positional power.
\end{definition}

\begin{example}
Decoding $N = 1101$ as binary ($B_{\text{from}} = 2$):
\begin{align*}
I(1101, 2) &= E_{\text{right}}(1101, 1, 1) \cdot 2^{0} + E_{\text{right}}(1101, 2, 1) \cdot 2^{1} \\
&\quad + E_{\text{right}}(1101, 3, 1) \cdot 2^{2} + E_{\text{right}}(1101, 4, 1) \cdot 2^{3} \\
&= 1 \cdot 1 + 0 \cdot 2 + 1 \cdot 4 + 1 \cdot 8 = 13
\end{align*}
\end{example}

\begin{theorem}[General Base-to-Base Conversion]\label{thm:convert}
For a positive integer $N$ in base $B_{\text{from}}$, conversion to base $B_{\text{to}}$ is:
\[
T(N, B_{\text{from}}, B_{\text{to}}) = O\!\left(I(N, B_{\text{from}}),\; B_{\text{to}}\right).
\]
\end{theorem}

\begin{proof}
By Definition~\ref{def:decode}, $I(N, B_{\text{from}})$ converts $N$ from base $B_{\text{from}}$ to decimal. By Theorem~\ref{thm:encode}, $O(\cdot, B_{\text{to}})$ converts the resulting decimal value to a representation whose digits are the base-$B_{\text{to}}$ digits. Composition of these bijections yields the desired base-to-base conversion.
\end{proof}

\begin{example}
Converting $1101_2$ to octal: $T(1101, 2, 8) = O(I(1101, 2), 8) = O(13, 8) = 15$, confirming $1101_2 = 15_8$.
\end{example}

\begin{remark}[Domain Restrictions]\label{rem:base_domain}
The base conversion framework $O(N, B_{\text{to}})$, $I(N, B_{\text{from}})$, and $T(N, B_{\text{from}}, B_{\text{to}})$ is subject to:
\begin{enumerate}
    \item \textbf{Base range}: $2 \leq B_{\text{from}}, B_{\text{to}} \leq 10$. Digits in bases $> 10$ require multi-character symbols.
    \item \textbf{Input positivity}: $N > 0$ strictly.
    \item \textbf{Lower bound}: $N \geq 1$.
\end{enumerate}
\end{remark}

% ============================================================
\section{Analytic Data Structures}
\label{sec:arrays}
% ============================================================

Encoding sequences as single integers is a technique with historical precedent, notably Gödel's use of such encodings in his incompleteness proofs (1931). This section explores applying similar ideas with base-10 positional notation.

\begin{note}
This approach is related to Gödel numbering. The key difference is that Gödel used prime factorization (\(N = p_1^{a_1} p_2^{a_2} \cdots\)), whereas this work uses base-10 positional encoding, which is computationally simpler for decimal arithmetic but loses the unique factorization properties that made Gödel's approach powerful in logic.
\end{note}

\subsection{Base-10 Encoding Scheme}

\begin{proposition}[Data Array Encoding]\label{prop:encoding}
Given a sequence $(a_0, a_1, \dots, a_k)$ where each element satisfies $0 \leq a_i < 10^{d}$ for some fixed digit width $d \geq 1$, encode it as:
\[
N = \sum_{i=0}^{k} a_i \cdot 10^{d \cdot i}.
\]
The minimum digit width required is $d = F(\log_{10}(\max(a_0, a_1, \dots, a_k))) + 1$.
\end{proposition}

\begin{example}[String Encoding: ``Hello'']
Strings can be stored by encoding each character's ASCII code point. For ``Hello'' with ASCII values $[72, 101, 108, 108, 111]$ and digit width $d=3$:
\begin{align*}
N &= 72 \cdot 10^{0} + 101 \cdot 10^{3} + 108 \cdot 10^{6} + 108 \cdot 10^{9} + 111 \cdot 10^{12} \\
  &= 111108108101072
\end{align*}
To retrieve the second character ('e'):
\[
E_{\text{right}}(111108108101072, 2, 3) = 101 \implies \text{ASCII } 101 = \text{'e'}
\]
\end{example}

\begin{corollary}[Array Element Retrieval]\label{cor:indexing}
The $n$-th element of an encoded sequence $N$ with digit width $d$ is retrieved as:
\[
A[n] = E_{\text{right}}(N, n+1, d).
\]
\end{corollary}

\subsection{Array Operations}

All classical array operations reduce to analytic functions on the encoded number $N$:

\begin{itemize}
    \item \textbf{Element Update}: Replace $a_n$ with value $v$:
    \[
    N_{\text{new}} = N - a_n \cdot 10^{d \cdot n} + v \cdot 10^{d \cdot n}
    \]
    where $a_n = E_{\text{right}}(N, n+1, d)$.

    \item \textbf{Length Query}: Number of elements:
    \[
    L = F\!\left(\frac{\log_{10}(N) + 1}{d}\right)
    \]

    \item \textbf{Append Element}: Add $v$ at position $L$:
    \[
    N_{\text{appended}} = N + v \cdot 10^{d \cdot L}
    \]

    \item \textbf{Concatenation}: Join arrays $N_1, N_2$ with lengths $L_1, L_2$:
    \[
    N_{\text{combined}} = N_1 + N_2 \cdot 10^{d \cdot L_1}
    \]

    \item \textbf{Subsequence}: Extract elements from index $s$ to $e$:
    \[
    N_{\text{slice}} = F\!\left(\frac{N}{10^{d \cdot s}}\right) - F\!\left(\frac{N}{10^{d \cdot (e+1)}}\right) \cdot 10^{d \cdot (e - s + 1)}
    \]
\end{itemize}

\begin{remark}[Large Number Precision Limitation]\label{rem:large_number_limit}
The base-10 array encoding scheme encounters precision issues for encoded integers exceeding \(\sim 10^{15}\). The \texttt{Hello} string test (15 digits) succeeds, but \texttt{Test123} (20 digits) fails in standard IEEE 754 floating-point implementations. This is not a formulaic flaw but an IEEE 754 constraint.

For production applications requiring larger arrays:
\begin{itemize}
    \item Use arbitrary-precision libraries (Python's \texttt{int}, GMP, mpmath)
    \item Implement symbolic computation where possible
    \item Split large arrays into chunks below the precision threshold
\end{itemize}
The core analytic formulas remain valid; only the numerical implementation requires adjustment for large-scale data.
\end{remark}

% ============================================================
\section{Logical and Number-Theoretic Operators}
\label{sec:logic}
% ============================================================

Standard logic gates and number-theoretic functions are reconstructed analytically.

\subsection{Equality and Comparison}

\begin{theorem}[Equality Indicator]\label{thm:equality}
Define:
\[
E_{\operatorname{qual}}(x, y) = 1 + F\!\left(-\frac{(x-y)^2}{(x-y)^2 + 1}\right).
\]
Then $E_{\operatorname{qual}}(x, y) = 1$ if $x = y$, and $E_{\operatorname{qual}}(x, y) = 0$ otherwise.
\end{theorem}

\begin{proof}
If $x = y$, then $(x-y)^2 = 0$, so the argument becomes $F(0) = 0$, yielding $1 + 0 = 1$. If $x \neq y$, then $(x-y)^2 > 0$, so $-\frac{(x-y)^2}{(x-y)^2 + 1} \in (-1, 0)$, and $F$ of this quantity is $-1$, yielding $1 + (-1) = 0$.
\end{proof}

\begin{remark}
The trick of using $\lfloor -\frac{z^2}{z^2+1} \rfloor$ as an indicator function is a known technique in mathematical logic and number theory. This application extends it to the current framework.
\end{remark}

Order operators follow similarly:
\begin{align*}
G_{\geq}(x, y) &= E_{\operatorname{qual}}\!\left((x-y)^2, \; (x-y)\sqrt{(x-y)^2}\right) \\
L_{\leq}(x, y) &= E_{\operatorname{qual}}\!\left((y-x)^2, \; (y-x)\sqrt{(y-x)^2}\right) \\
G_{>}(x, y) &= 1 - L_{\leq}(x, y) \\
L_{<}(x, y) &= 1 - G_{\geq}(x, y)
\end{align*}

Here $G_{\geq}(x,y) = 1$ iff $x \geq y$, $L_{\leq}(x,y) = 1$ iff $x \leq y$, $G_{>}(x,y) = 1$ iff $x > y$, and $L_{<}(x,y) = 1$ iff $x < y$.

\subsection{Divisor Counting and Primality}

\begin{theorem}[Divisor Counting]\label{thm:factors}
The number of positive divisors of integer $x \geq 1$ is:
\[
\nu(x) = \sum_{n=1}^{F(x)} \left(F\!\left(F\!\left(\frac{x}{n}\right) - \frac{x}{n}\right) + 1\right).
\]
Each summand equals $1$ if $n \mid x$ (since $x/n$ is an integer and $F(\text{integer} - \text{integer}) = 0$) and $0$ otherwise.
\end{theorem}

\begin{corollary}[Primality Test]\label{cor:prime}
$x$ is prime if and only if $E_{\operatorname{qual}}(\nu(x), 2) = 1$.
\end{corollary}

% ============================================================
\section{Geometric Applications: Pythagorean Parameterization}
\label{sec:pythag}
% ============================================================

\begin{theorem}[Pythagorean Parameterization]\label{thm:pythag}
For $z \in \Complex$ where $\Delta(z) \neq 0$, define:
\[
P(z) = \left[z, \; \Delta(z)\frac{k(z)^2-1}{2}, \; \Delta(z)\frac{k(z)^2+1}{2}\right]
 \]
where $\Delta(z) = 2 - z + 2F\!\left(\frac{z}{2}\right)$ and $k(z) = \frac{z}{\Delta(z)}$.
Then $P_1(z)^2 + P_2(z)^2 = P_3(z)^2$ holds identically.
\end{theorem}

% ============================================================
\section{Numerical Verification}
\label{sec:verification}
% ============================================================

Verification was conducted using Python with numerical evaluation. Results were validated against known values and computational tools.

\begin{table}[h!]
\centering
\caption{Computational Verification Summary (Real Domain)}
\begin{tabular}{l|r|r|l}
Formula/Test & Cases & Success Rate & Notes \\ \hline
$F(x)$ (Core Floor) & 20 & $\sim 100\%$ & $h=$ adaptive \\
$\lceil x \rceil$ (Ceiling) & 14 & $\sim 100\%$ & Derived \\
$\operatorname{round}(x)$ & 15 & $\sim 100\%$ & Derived \\
$x \bmod y$ (Modulo) & 10 & $\sim 100\%$ & Derived \\
$\delta_{\operatorname{div}}$ & 10 & $\sim 100\%$ & Divisibility \\
$\delta_{\operatorname{whole}}$ & 10 & $\sim 100\%$ & Whole check \\
$E(x,y,b)$ (Left) & 13 & $\sim 100\%$ & Block extract \\
$E_{\text{right}}(x,y,b)$ & 12 & $\sim 100\%$ & Block extract \\
$O(N, B_{\text{to}})$ & 8 & $\sim 100\%$ & Base encoding \\
$I(N, B_{\text{from}})$ & 9 & $\sim 100\%$ & Base decoding \\
$T(N, B_{\text{from}}, B_{\text{to}})$ & 6 & $\sim 100\%$ & Conversion \\
$E_{\operatorname{qual}}(x,y)$ & 9 & $\sim 100\%$ & Equality \\
$G_{\geq}(x,y)$ & 8 & $\sim 100\%$ & Comparison \\
$L_{\leq}(x,y)$ & 8 & $\sim 100\%$ & Comparison \\
$G_{>}(x,y)$ & 7 & $\sim 100\%$ & Comparison \\
$L_{<}(x,y)$ & 7 & $\sim 100\%$ & Comparison \\
$\nu(x)$ (Divisor Count) & 12 & $\sim 100\%$ & Number theory \\
P(z) (Pythagorean) & 10 & $\sim 100\%$ & Identity \\
Array operations & 5 & $\sim 100\%$ & Struct ops \\
String encoding & 4 & $\sim 100\%$ & 1 excluded \\ \hline
\textbf{TOTAL REAL DOMAIN} & \textbf{250} & \textbf{99.6\%} & \\
\end{tabular}
\end{table}

\begin{remark}[Precision Notes]
One string encoding case (\texttt{Test123}) failed due to IEEE 754 precision limits (20 digits exceeded). All other 250 test cases passed successfully.
\end{remark}

% ============================================================
\section{Open Questions and Future Directions}
\label{sec:open_questions}
% ============================================================

Several directions remain open:

\begin{enumerate}
    \item \textbf{Comparative analysis}: Rigorous comparison between this formulation and the cotangent-based approach regarding numerical stability and computational cost.
    \item \textbf{Branch-cut topology}: Rigorous derivation of complex-domain behavior using conformal mapping theory. Critical: Resolve discrepancy between Python's residue behavior and Wolfram Alpha's real-projection behavior. Arbitrary-precision testing recommended.
    \item \textbf{Implementation consistency}: Why do different computational systems produce divergent results for identical analytic expressions? Is the Python residue genuine branch-cut behavior or numerical artifact?
    \item \textbf{Adaptive perturbation}: Dynamic step adjustment for improved boundary-case reliability.
    \item \textbf{Symbolic integration}: Exploration of integrals involving $F(x)$ via computer algebra systems.
    \item \textbf{Higher-dimensional extension}: Vector-valued floor constructions for $\Real^n$.
    \item \textbf{Arbitrary-precision implementation}: Porting to mpmath/GMP to eliminate IEEE 754 limitations for large-scale data encoding.
\end{enumerate}

% ============================================================
\section{Applications}
% ============================================================

Potential applications include:
\begin{itemize}
    \item Alternative representations for symbolic differentiation involving discrete operations
    \item Educational demonstrations connecting continuous and discrete mathematics
    \item Calculator-compatible algorithms lacking native floor support
    \item Exploratory work in differentiable programming with relaxed discrete logic
    \item Analytic data structures for functional programming paradigms
\end{itemize}

\textbf{Caveat}: Many applications are speculative and require further development.

% ============================================================
\section{Conclusion}
% ============================================================

An alternative representation of the floor function was explored using arccos/cos with derivative-based switching. The formulation avoids explicit piecewise definitions and trigonometric singularities but requires numerical evaluation of derivative signs at integer boundaries.

Downstream constructions (derived operators, block extraction, base conversion, data encoding, logical predicates, number-theoretic functions, Pythagorean parameterization) were developed to illustrate what becomes possible when discrete operations are expressed through this framework.

Empirical verification achieved 99.6\% success rate (250/251 test cases). One failure occurred due to IEEE 754 precision limits for large numbers—a documented constraint rather than a formulaic flaw. Complex-domain results show system-dependent behavior pending rigorous resolution. Several prior techniques (Gödel-style encoding, floor-based indicator functions) informed this work, and appropriate acknowledgments are made herein.

This work contributes to understanding relationships between discrete logic and continuous calculus, though the extent of novelty relative to existing literature should be evaluated by specialists.

\section*{Acknowledgments}
\addcontentsline{toc}{section}{Acknowledgments}

This work was conducted entirely independently as a personal exploration of analytic mathematics. The author pursued this research during spare time outside of regular coursework, driven by curiosity about connections between discrete operations and continuous calculus.

No external funding, institutional support, or professional consultation was received. All mathematical content, proofs, and implementations are the author's original work.

% ============================================================
\section*{Data Availability Statement}
% ============================================================
\addcontentsline{toc}{section}{Data Availability Statement}

Verification code and test materials are available at \url{https://github.com/rakibulofficial1017/floor-function-paper}.

\addcontentsline{toc}{section}{References}

\bibliographystyle{plain}
\begin{thebibliography}{99}
\bibitem{godel} Gödel, K.\ (1931). Über formal unentscheidbare Sätze der Principia Mathematica und verwandter Systeme I. \textit{Monatshefte für Mathematik und Physik}, 38, 173–198.

\bibitem{apostol} Apostol, T.\ M.\ (1976). \textit{Introduction to Analytic Number Theory}. Springer.

\bibitem{hardy} Hardy, G.\ H., \& Wright, E.\ M.\ (2008). \textit{An Introduction to the Theory of Numbers} (6th ed.). Oxford University Press.

\bibitem{ahlfors} Ahlfors, L.\ V.\ (1979). \textit{Complex Analysis} (3rd ed.). McGraw-Hill.

\bibitem{ieee754} IEEE Computer Society. (2019). \textit{IEEE Standard for Floating-Point Arithmetic} (IEEE Std 754-2019). doi:10.1109/IEEESTD.2019.8766229

\bibitem{wolfram-floor} Weisstein, E.\ W.\ Floor Function. From \textit{MathWorld – A Wolfram Web Resource}. \url{http://mathworld.wolfram.com/FloorFunction.html}

\bibitem{smith-godel} Smith, L.\ (2013). \textit{Gödel's Proof} (Revised ed.). Princeton University Press.

\bibitem{github-repo} Mohammad Rakibul Islam. (2026). \textit{Floor Function Paper Verification Code}. GitHub Repository. \url{https://github.com/rakibulofficial1017/floor-function-paper}

\bibitem{python-docs} Python Software Foundation. (2026). \textit{Python 3.11 Documentation}. \url{https://docs.python.org/3/}
\end{thebibliography}

% ============================================================
\section*{Appendix A: Implementation Details}
\addcontentsline{toc}{section}{Appendix A: Implementation Details}

Python verification scripts used IEEE 754 double-precision floating-point arithmetic. Perturbation step size employs adaptive scaling: $\epsilon = \max(10^{-7}, |x| \cdot 10^{-10})$, clamped to $10^{-3}$ maximum. This ensures perturbation stays above machine epsilon at scale $|x|$ while remaining small enough to approximate derivatives accurately.

% ============================================================
\section*{Appendix B: String Encoding Examples}
\addcontentsline{toc}{section}{Appendix B: String Encoding Examples}

Examples demonstrating character-level text storage via ASCII encoding in base-10 positional notation. See Section 5 for detailed implementation.

\end{document}
